Задачи машинного обучения с учителем, как правило, состоят в восстановлении зависимости между парами (признаковое
описание, целевая переменная) по данным, доступным нам для анализа.
Алгоритмы машинного обучения (learning algorithm), со многими из которых вы уже успели познакомиться, позволяют
построить модель, аппроксимирующую эту зависимость.
Но как понять, насколько качественной получилась аппроксимация?

Почти наверняка наша модель будет ошибаться на некоторых объектах: будь она даже идеальной, шум или выбросы в
тестовых данных всё испортят.
При этом разные модели будут ошибаться на разных объектах и в разной степени.
Задача специалиста по машинному обучению -- подобрать подходящий критерий, который позволит сравнивать различные
модели.

Важно: качество модели нельзя оценивать на обучающей выборке.
Как минимум, это стоит делать на отложенной (тестовой) выборке, но если вам это позволяют время и вычислительные
ресурсы, стоит прибегнуть и к более надёжным способам проверки — например, \textit{кросс-валидации}

\underline{Выбор метрик в реальных задачах}
Возможно, вы уже участвовали в соревнованиях по анализу данных.
На таких соревнованиях метрику (критерий качества модели) организатор выбирает за вас, и она, как правило, довольно
понятным образом связана с результатами предсказаний.
Но на практике всё бывает намного сложнее.
Варианты наших действий

\begin{itemize}
    \item Решить, сколько коробок с бананами нужно завтра привезти в конкретный магазин -- чтобы предложение
    соответствовало спросу, и не пришлось выбрасывать излишки;
    \item увеличить счастье пользователя от работы с сервисом, чтобы он стал лояльным и приносил стабильный
    прогнозируемый доход;
    \item решить, нужно ли направить человека на дополнительное обследование.
\end{itemize}

В каждом конкретном случае может возникать целая иерархия метрик.
Представим, например, что речь идёт о стриминговом музыкальном сервисе, пользователей которого мы решили порадовать
сгенерированными самодельной нейросетью треками -- не защищёнными авторским правом, а потому совершенно бесплатными.

Иерархия метрик

\begin{enumerate}
    \item Самый верхний уровень: будущий доход сервиса -- невозможно измерить в моменте, сложным образом зависит
    от совокупности всех наших усилий.
    \item Медианная длина сессии, возможно, служащая оценкой радости пользователей, которая, как мы надеемся,
    повлияет на их желание продолжать платить за подписку -- её нам придётся измерять в продакшене, ведь нас
    интересует реакция настоящих пользователей на новшество.
    \item Доля удовлетворённых качеством сгенерированной музыки асессоров, на которых мы потестируем её до того,
    как выставить на суд пользователей.
    \item Функция потерь, на которую мы будем обучать генеративную сеть.
\end{enumerate}
На этом примере мы можем заметить сразу несколько общих закономерностей.
Во-первых, метрики бывают offline и online.
Online-метрики вычисляются по данным, собираемым с работающей системы (например, медианная длина сессии).
Offline-метрики могут быть измерены до введения модели в эксплуатацию, например, по историческим данным или с
привлечением специальных людей, асессоров.

Последнее часто применяется, когда метрика -- это реакция живого человека: скажем, так поступают поисковые
компании, которые предлагают людям оценить качество ранжирования экспериментальной системы ещё до того, как рядовые
пользователи увидят эти результаты в обычном порядке.
На самом же нижнем этаже иерархии лежат оптимизируемые в ходе обучения функции потерь.

В данном разделе нас будут интересовать offline-метрики, которые могут быть измерены без привлечения людей

\subsection{Определение метрики в ML}

Метрика в машинном обучении -- это показатель, который количественно характеризует качество и свойство
компьютерного алгоритма

\subsubsection{Свойства, сравнение с функцией потерь}

Свойства метрик могут включать простоту интерпретации, универсальность, учёт дисбаланса классов и другие.
Важно помнить, что ни одна метрика не идеальна и профессиональный подход требует рассмотрения комбинации метрик,
а также их соотнесения с бизнес-требованиями

\underline{Функция потерь не есть метрика качества}
Как мы узнали ранее, методы обучения реализуют разные подходы

\begin{itemize}
    \item Обучение на основе прироста информации (как в деревьях решений);
    \item обучение на основе сходства (как в методах ближайших соседей);
    \item обучение на основе вероятностной модели данных (например, максимизацией правдоподобия);
    \item обучение на основе ошибок (минимизация эмпирического риска).
\end{itemize}

И в рамках обучения на основе минимизации ошибок мы уже отвечали на вопрос: как можно штрафовать модель за
предсказание на обучающем объекте.

Во время сведения задачи о построении решающего правила к задаче численной оптимизации, мы вводили понятие функции
потерь и, обычно, объявляли целевой функцией сумму потерь от предсказаний на всех объектах обучающей выборки.

Важно понимать разницу между функцией потерь и метрикой качества.
Её можно сформулировать следующим образом:

\begin{itemize}
    \item Функция потерь возникает в тот момент, когда мы сводим задачу построения модели к задаче оптимизации.
    Обычно требуется, чтобы она обладала хорошими свойствами (например, дифференцируемостью).
    \item Метрика — внешний, объективный критерий качества, обычно зависящий не от параметров модели, а только от
    предсказанных меток.
\end{itemize}

В некоторых случаях метрика может совпадать с функцией потерь.
Например, в задаче регрессии MSE играют роль как функции потерь, так и метрики.
Но, скажем, в задаче бинарной классификации они почти всегда различаются: в качестве функции потерь может выступать
кросс-энтропия, а в качестве метрики -- число верно угаданных меток (accuracy).
Отметим, что в последнем примере у них различные аргументы: на вход кросс-энтропии нужно подавать logit, а на вход
accuracy — предсказанные метки (то есть по сути argmax logit).

Начнём с самой простой разновидности классификации -- бинарной, а затем постепенно будем наращивать сложность.

Напомним постановку задачи бинарной классификации: нам нужно по обучающей выборке
${(x_i,y_i)}_{i=1}^N$, где $y_i \in {0,1}$ построить модель, которая по объекту $x$ предсказывает метку класса
$f(x) \in {0,1}$.

Первый критерий качества, который приходит в голову, -- accuracy, то есть доля объектов с правильным предсказанием
класса:
\[ Accuracy(y,y_{pred}) = \frac{1}{N} \sum_{i=1}^N \mathbb{I} [y_i = f(x_i)] \]

Или же сопряжённая ей метрика -- доля ошибочных классификаций (error rate): Error rate = 1 − Accuracy

Познакомившись чуть внимательнее с этой метрикой, можно заметить, что у неё есть несколько недостатков

\begin{itemize}
    \item Она не учитывает дисбаланс классов.
    Например, в задаче диагностики редких заболеваний классификатор, предсказывающий всем пациентам отсутствие
    болезни будет иметь достаточно высокую accuracy просто потому, что больных людей в выборке намного меньше;
    \item она также не учитывает цену ошибки на объектах разных классов.
    Для примера снова можно привести задачу медицинской диагностики: если ошибочный положительный диагноз для
    здорового больного обернётся лишь ещё одним обследованием, то ошибочно отрицательный вердикт может повлечь
    роковые последствия.
\end{itemize}

\subsection{Матрица ошибок (бинарная, многоклассовая)}

Исторически задача бинарной классификации -- это задача об обнаружении чего-то редкого в большом потоке объектов,
например, поиск человека, больного туберкулёзом, по флюорографии.
Или задача признания пятна на экране приёмника радиолокационной станции бомбардировщиком, представляющем угрозу
охраняемому объекту (в противовес стае гусей).

Для задачи классификации с $C$ классами, \textbf{матрица ошибок} (confusion matrix) $\mathbf{M}$ размера $C \times C$
определяется как:

\[ \mathbf{M} = [m_{ij}]_{i,j=1}^{C} \]
где:
\begin{itemize}
    \item $m_{ij}$ --- количество объектов класса $i$, отнесённых алгоритмом к классу $j$
    \item $i$ --- индекс истинного класса (строка)
    \item $j$ --- индекс предсказанного класса (столбец)
\end{itemize}

\underline{Пример для 3 классов}

\begin{table}[h]
    \centering
    \begin{tabular}{|c|c|c|c|c|}
        \hline
        \multirow{2}{*}{\textbf{Истина $\backslash$ Предсказание}} & \multicolumn{3}{c|}{\textbf{Предсказанный класс}}
        & \multirow{2}{*}{\textbf{Всего}} \\ \cline{2-4}
        & Класс 1           & Класс 2           & Класс 3           &       \\ \hline
        \textbf{Класс 1} & $m_{11}$ (TP$_1$) & $m_{12}$          & $m_{13}$          & $N_1$ \\ \hline
        \textbf{Класс 2} & $m_{21}$          & $m_{22}$ (TP$_2$) & $m_{23}$          & $N_2$ \\ \hline
        \textbf{Класс 3} & $m_{31}$          & $m_{32}$          & $m_{33}$ (TP$_3$) & $N_3$ \\ \hline
        \textbf{Всего}   & $P_1$             & $P_2$             & $P_3$             & $N$   \\ \hline
    \end{tabular}
    \caption{Матрица ошибок для трёх классов}
    \label{tab:triple matrix}
\end{table}

\underline{Вычисление основных показателей для класса $c$}
Для каждого класса $c$ можно вычислить:
\begin{align*}
    \text{TP}_c &= m_{cc} \quad \text{(True Positives для класса $c$).} \\
    \text{FP}_c &= \sum_{i \neq c} m_{ic} \quad \text{(False Positives для класса $c$).} \\
    \text{FN}_c &= \sum_{j \neq c} m_{cj} \quad \text{(False Negatives для класса $c$).} \\
    \text{TN}_c &= \sum_{i \neq c} \sum_{j \neq c} m_{ij} \quad \text{(True Negatives для класса $c$)}
\end{align*}

\subsubsection{Типы ошибок}

Класс, который представляет для нас интерес, называется «положительным», а оставшийся — «отрицательным».
Заметим, что для каждого объекта в выборке возможно 4 ситуации

\begin{itemize}
    \item Мы предсказали \textit{положительную} метку и \textit{угадали}.
    Будет относить такие объекты к true positive (TP) группе.
    True — потому что предсказали мы правильно, а positive — потому что предсказали положительную метку;
    \item мы предсказали \textit{положительную} метку, но \textit{ошиблись} в своём предсказании -- false
    positive (FP).
    False, потому что предсказание было неправильным;
    \item мы предсказали \textit{отрицательную} метку и \textit{угадали} — true negative (TN);
    \item и наконец, мы предсказали \textit{отрицательную} метку, но \textit{ошиблись} -- false negative (FN).
    Для удобства все эти 4 числа изображают в виде таблицы, которую называют confusion matrix (матрицей ошибок):
\end{itemize}

В многоклассовой задаче каждый класс рассматривается как \("\)положительный\("\), а все остальные вместе как
\("\)отрицательный\("\):

\begin{itemize}
    \item \textbf{True Positive для класса $c$}: Объект класса $c$ правильно классифицирован как $c$
    \item \textbf{False Positive для класса $c$}: Объект не из класса $c$ ошибочно классифицирован как $c$
    \item \textbf{False Negative для класса $c$}: Объект класса $c$ ошибочно классифицирован как другой класс
    \item \textbf{True Negative для класса $c$}: Объект не из класса $c$ правильно классифицирован как не $c$
\end{itemize}

\subsubsection{Правило присвоения меток}

\begin{enumerate}
    \item \textbf{Пороговое правило для вероятностей:}
    Для бинарной классификации с порогом $\tau$:
    \[ \hat{y}_i =
    \begin{cases}
        1, & \text{если } \hat{p}_i \geq \tau \\
        0, & \text{если } \hat{p}_i < \tau
    \end{cases} \]
    
    \item \textbf{Argmax для многоклассовой:}
    \[ \hat{y}_i = \arg\max_{c \in \{1,\dots,C\}} \hat{p}_{ic} \]
    где $\hat{p}_{ic}$ --- предсказанная вероятность принадлежности к классу $c$
    
    \item \textbf{Однопроходное правило:}
    \[ \hat{y}_i =
    \begin{cases}
        1, & \text{если } f(\bm{x}_i) \geq 0 \\
        0, & \text{если } f(\bm{x}_i) < 0
    \end{cases} \]
    для классификатора $f: \mathbb{R}^d \to \mathbb{R}$
\end{enumerate}

\underline{Стратегии для несбалансированных данных}

\begin{itemize}
    \item \textbf{Настройка порога:} смещение порога в сторону меньшего класса
    \[ \tau_{\text{adjusted}} = \frac{N_{\text{negative}}}{N_{\text{positive}} + N_{\text{negative}}} \]
    
    \item \textbf{Разные пороги для разных классов:}
    \[ \hat{y}_i = c \quad \text{если} \quad \frac{\hat{p}_{ic}}{\tau_c} > \frac{\hat{p}_{ij}}{\tau_j} \ \forall j \neq
    c \]
    
    \item \textbf{Правило с учётом стоимости ошибок:}
    \[ \hat{y}_i = \arg\min_{c} \sum_{j=1}^C L(j,c) \cdot \hat{p}_{ij} \]
    где $L(j,c)$ --- стоимость ошибки классификации класса $j$ как класса $c$
\end{itemize}

\subsection{Accuracy, balanced accuracy, precision, recall, F1-score}
Accuracy был выше.

\begin{definition}[Balanced Accuracy]
    Для бинарной классификации:
    \[ \text{Balanced Accuracy} = \frac{1}{2} \left( \frac{TP}{TP+FN} + \frac{TN}{TN+FP} \right) =
    \frac{\text{Sensitivity} + \text{Specificity}}{2} \]
    
    Для многоклассовой классификации:
    \[ \text{Balanced Accuracy} = \frac{1}{C} \sum_{c=1}^C \frac{TP_c}{TP_c + FN_c} \]
\end{definition}

\textbf{Преимущества}
\begin{itemize}
    \item Устойчива к дисбалансу классов.
    \item Даёт равный вес каждому классу
    \item Лучше отражает качество на несбалансированных данных
\end{itemize}


В случае асимметрии классов, можно использовать метрики, которые не учитывают TN и ориентируются на TP\@.
Если мы рассмотрим долю правильно предсказанных положительных объектов среди всех объектов, предсказанных
положительным классом, то мы получим метрику, которая называется точностью (precision)
\[ Precision = \frac{TP}{TP+FP} \]
Интуитивно метрика показывает долю релевантных документов среди всех найденных классификатором.
Чем меньше ложноположительных срабатываний будет допускать модель, тем больше будет её Precision.

Если же мы рассмотрим долю правильно найденных положительных объектов среди всех объектов положительного класса, то
мы получим метрику, которая называется полнотой (recall)
\[ Recall = \frac{TP}{TP+FN} \]
Интуитивно метрика показывает долю найденных документов из всех релевантных.
Чем меньше ложно отрицательных срабатываний, тем выше recall модели.

Например, в задаче предсказания злокачественности опухоли точность показывает, сколько из определённых нами как
злокачественные, действительно таковые, а полнота -- какую долю злокачественных опухолей нам удалось выявить.

\underline{$F_1$-мера}

Как мы уже отмечали ранее, модели очень удобно сравнивать, когда их качество выражено одним числом.
В случае пары Precision-Recall существует популярный способ скомпоновать их в одну метрику -- взять их среднее
гармоническое.
Данный показатель эффективности исторически носит название $F_1$-меры ($F_1$-measure).
\[ F_1 = \frac{2}{\frac{1}{Recall} + \frac{1}{Precision}} \]

Стоит иметь в виду, что $F_1$-мера предполагает одинаковую важность Precision и Recall, если одна из этих метрик для
вас приоритетнее, то можно воспользоваться $F_\beta$ мерой:
\[ F_1 = (\beta^2 + 1) \frac{Recall  \dot Precision}{Recall + \beta^2 Precision} \]

\subsubsection{Sensitivity, specificity и другие}

Sensitivity = recall.
В медицинских приложениях:
\begin{itemize}
    \item \textbf{Sensitivity (Чувствительность):} Вероятность того, что тест будет положительным при наличии заболевания
    \[ \text{Sensitivity} = P(\text{Тест} + \mid \text{Болезнь} +) = \frac{TP}{TP+FN} \]
    
    \item \textbf{Specificity (Специфичность):} Вероятность того, что тест будет отрицательным при отсутствии
    заболевания или доля найденных отрицательных в общём объёме отрицательных
    \[ \text{Specificity} = P(\text{Тест} - \mid \text{Болезнь} -) = \frac{TN}{TN+FP} \]
\end{itemize}

\underline{Расширенные метрики}

\begin{enumerate}
    \item \textbf{Matthews Correlation Coefficient (MCC):}
    \[ \text{MCC} = \frac{TP \cdot TN - FP \cdot FN}{\sqrt{(TP+FP)(TP+FN)(TN+FP)(TN+FN)}} \]
    Корреляция между истинными и предсказанными метками
    
    \item \textbf{Cohen's Kappa} Согласованность классификаций с учётом случайного угадывания
    \[ \kappa = \frac{p_o - p_e}{1 - p_e} \]
    где $p_o$ --- наблюдаемая точность, $p_e$ --- ожидаемая точность при случайном угадывании
\end{enumerate}

\subsubsection{Несбалансированная стоимость ошибок, примеры, как исправлять дисбаланс}

\underline{Матрица стоимости ошибок}

Для задачи с $C$ классами задаётся матрица стоимости $\mathbf{L} \in \mathbb{R}^{C \times C}$:
\[ \mathbf{L} = [L_{ij}]_{i,j=1}^C \]
Здесь $L_{ij}$ --- стоимость классификации объекта класса $i$ как класса $j$.

\textbf{Свойства:}
\begin{itemize}
    \item $L_{ii} = 0$ (правильная классификация бесплатна)
    \item $L_{ij} > 0$ для $i \neq j$ (ошибки имеют положительную стоимость)
    \item Часто $L_{ij} \neq L_{ji}$ (стоимость ошибок несимметрична)
\end{itemize}

\underline{Примеры из реальных задач}

\begin{table}[h]
    \centering
    \begin{tabular}{|p{0.3\textwidth}|p{0.3\textwidth}|p{0.3\textwidth}|}
        \hline
        \textbf{Задача}         & \textbf{Дорогая ошибка}          & \textbf{Менее дорогая ошибка}     \\ \hline
        Медицинская диагностика & Пропустить болезнь (FN)          & Ложная тревога (FP)               \\ \hline
        Спам-фильтрация         & Пропустить спам (FN)             & Пометить не-спам как спам (FP)    \\ \hline
        Кредитный scoring       & Выдать кредит неплательщику (FP) & Отказать плательщику (FN)         \\ \hline
        Автономное вождение     & Не заметить пешехода (FN)        & Ложное срабатывание тормозов (FP) \\ \hline
    \end{tabular}
    \caption{Примеры несбалансированной стоимости ошибок}
    \label{tab:examples}
\end{table}

\underline{Способы учёта различной стоимости ошибок}

\underline{1. Взвешенная функция потерь}

Модификация функции потерь с весами:
\[ \mathcal{L}_{\text{weighted}}(\bm{w}) = \sum_{i=1}^n \alpha_{y_i} \cdot \ell(y_i, f(\bm{x}_i; \bm{w})) \]
Здесь $\alpha_c$ --- вес для класса $c$, обычно $\alpha_c \propto \frac{1}{N_c}$.

\underline{2. Изменение порога принятия решения}

Для бинарной классификации с разными стоимостями $C_{FP}$ и $C_{FN}$ оптимальный порог:
\[ \tau^* = \frac{C_{FP}}{C_{FP} + C_{FN}} \]
При условии калиброванных вероятностей.

\underline{3. Cost-sensitive learning}

Модификация алгоритмов для минимизации общей стоимости:
\[ \hat{\bm{w}} = \arg\min_{\bm{w}} \sum_{i=1}^n \sum_{j=1}^C L(y_i, j) \cdot \mathbb{I}[\hat{y}_i = j] \]

\underline{4. Методы выборки (sampling)}

\begin{itemize}
    \item \textbf{Oversampling:} дублирование примеров миноритарного класса
    \item \textbf{Undersampling:} удаление примеров мажоритарного класса
    \item \textbf{SMOTE (Synthetic Minority Oversampling):} генерация синтетических примеров миноритарного класса
\end{itemize}

\underline{Как исправлять дисбаланс классов}

\underline{На уровне данных}

\begin{enumerate}
    \item \textbf{Стратифицированная выборка:} сохранение пропорций классов при разбиении.
    \item \textbf{Взвешенная случайная выборка:} выборка с вероятностями, обратно пропорциональными размеру класса.
    \item \textbf{Генеративные методы:}
    \begin{itemize}
        \item SMOTE: создание синтетических примеров миноритарного класса
        \item ADASYN: адаптивный вариант SMOTE
        \item генеративно-состязательные сети (GAN)
    \end{itemize}
\end{enumerate}

\underline{На уровне алгоритма}

\begin{enumerate}
    \item \textbf{Взвешенная функция потерь:} присвоение больших весов ошибкам на миноритарном классе
    \item \textbf{Cost-sensitive learning:} явное задание матрицы стоимости ошибок
    \item \textbf{Ансамблирование:} использование методов вроде EasyEnsemble или BalanceCascade
\end{enumerate}

\underline{На уровне оценки}

\begin{enumerate}
    \item Использование метрик, устойчивых к дисбалансу:
    \begin{itemize}
        \item Balanced Accuracy
        \item F1-score, F$\beta$-score
        \item MCC, Cohen's Kappa
        \item AUC-ROC, AUC-PR
    \end{itemize}
    
    \item Стратифицированная кросс-валидация
\end{enumerate}

\underline{Практические рекомендации}

\begin{table}[h]
    \centering
    \begin{tabular}{|p{0.3\textwidth}|p{0.3\textwidth}|p{0.3\textwidth}|}
        \hline
        \textbf{Соотношение классов}      & \textbf{Рекомендуемый подход}  & \textbf{Метрики}            \\ \hline
        Небольшой дисбаланс (1:3)         & Взвешенная функция потерь      & F1-score, Balanced Accuracy \\ \hline
        Средний дисбаланс (1:10)          & SMOTE + взвешенная loss        & AUC-PR, MCC                 \\ \hline
        Сильный дисбаланс (1:1000+) & One-class classification       & Precision, Specificity      \\ \hline
    \end{tabular}
    \caption{Рекомендации по работе с различными уровнями дисбаланса}
    \label{tab:disbalance recs}
\end{table}

\subsection{ROC AUC}

Пусть мы хотим учитывать ошибки на объектах обоих классов.
При уменьшении порога отсечения мы будем находить (правильно предсказывать) всё большее число положительных
объектов, но также и неправильно предсказывать положительную метку на всё большем числе отрицательных объектов.
Естественным, кажется, ввести две метрики TPR и FPR:

TPR (true positive rate) -- это полнота, доля положительных объектов, правильно предсказанных положительными:
\[ TPR=\frac{TP}{P} = \frac{TP}{TP+FN} \]
FPR (false positive rate) -- это доля отрицательных объектов, неправильно предсказанных положительными:
\[ FPR = \frac{FP}{N} = \frac{FP}{FP+TN} \]

Обе эти величины растут при уменьшении порога.
Кривая в осях TPR/FPR, которая получается при варьировании порога, исторически называется ROC-кривой (receiver
operating characteristics curve, сокращённо ROC curve).

Жёлтая и синяя кривые показывают распределение предсказаний классификатора на объектах положительного и
отрицательного классов соответственно.
То есть значения на оси X (на графике с двумя гауссианами) мы получаем из классификатора.

\begin{itemize}
    \item Если классификатор идеальный, -- две кривые разделимы по оси X, -- то на правом графике мы получаем
    ROC-кривую $(0,0)->(0,1)->(1,1)$, площадь под которой равна 1.
    \item Если классификатор случайный (предсказывает одинаковые метки положительным и отрицательным объектам), то
    мы получаем ROC-кривую $(0,0)->(1,1)$, площадь под которой равна 0.5.
\end{itemize}

\subsubsection{Свойства ROC и ROC-AUC}
Чем лучше классификатор разделяет два класса, тем больше площадь (\textit{area under curve}) под ROC-кривой -- и
мы можем использовать её в качестве метрики.
Эта метрика называется AUC и она работает благодаря следующему свойству ROC-кривой:

AUC равен доле пар объектов вида (объект класса 1, объект класса 0), которые алгоритм верно упорядочил, то есть
предсказание классификатора на первом объекте больше:

\begin{gather*}
    AUC = \frac{\sum\limits_{i=1}^N \sum\limits_{j=1}^N I[y_i < y_j] \cdot I'[f(x_i) < f(x_j)]}{\sum\limits_{i=1}^N
    \sum\limits_{j=1}^N I[y_i < y_j]} \\
    I'[f(x_i) < f(x_j)] =
    \begin{cases}
        0, & f(x_i) > f(x_j) \\
        0.5, & f(x_i) = f(x_j) \\
        1, & f(x_i) < f(x_j)
    \end{cases}\\
    I[y_i < y_j] =
    \begin{cases}
        0, & y_i \geq y_j \\
        1, & y_i < y_j
    \end{cases} \\
\end{gather*}

В каких случаях лучше отдать предпочтение этой метрике?
Рассмотрим следующую задачу: некоторый сотовый оператор хочет научиться предсказывать, будет ли клиент пользоваться
его услугами через месяц.
На первый взгляд кажется, что задача сводится к бинарной классификации с метками 1, если клиент останется с компанией и
0 если иначе.

Однако если копнуть глубже в процессы компании, то окажется, что такие метки практически бесполезны.
Компании скорее интересно упорядочить клиентов по вероятности прекращения обслуживания и в зависимости от этого
применять разные варианты удержания: кому-то прислать скидочный купон от партнёра, кому-то предложить скидку на
следующий месяц, а кому-то и новый тариф на особых условиях.

Таким образом, в любой задаче, где нам важна не метка сама по себе, а правильный порядок на объектах, имеет смысл
применять AUC\@.
Утверждение выше может вызывать у вас желание использовать AUC в качестве метрики в задачах ранжирования, но мы
призываем вас быть аккуратными.

\subsubsection{PR AUC}

\begin{definition}[PR кривая]
    Кривая Precision-Recall строится путём варьирования порога классификации $\tau$ и вычисления:
    \begin{itemize}
        \item Precision (точность) на оси Y
        \item Recall (полнота) на оси X
    \end{itemize}
    Для каждого порога $\tau$:
    \[ \text{Precision}(\tau) = \frac{TP(\tau)}{TP(\tau) + FP(\tau)}, \quad
        \text{Recall}(\tau) = \frac{TP(\tau)}{TP(\tau) + FN(\tau)} \]
\end{definition}

\underline{PR AUC vs ROC AUC}

\begin{table}[h]
    \centering
    \begin{tabular}{|p{0.4\textwidth}|p{0.25\textwidth}|p{0.25\textwidth}|} \hline
        \textbf{Характеристика.}       & \textbf{PR AUC}                  & \textbf{ROC AUC}      \\ \hline
        Чувствительность к дисбалансу & Высокая                          & Низкая                \\ \hline
        Интерпретация при дисбалансе  & Более информативна               & Может быть обманчивой \\ \hline
        Оси                           & Precision (Y) vs Recall (X)      & TPR (Y) vs FPR (X)    \\ \hline
        Базовая линия                 & Положительная доля               & Диагональ y = x       \\ \hline
        Подходит для                  & Сильно несбалансированных данных & Умеренного дисбаланса \\ \hline
    \end{tabular}
    \caption{Сравнение PR AUC и ROC AUC}
    \label{tab:pr roc comp}
\end{table}

\underline{Вычисление PR AUC}

\underline{1. Точное вычисление (по всем порогам)}

\begin{algorithmic}[1]
    \State Получить предсказанные вероятности $\hat{p}_i$ для всех объектов.
    \State Отсортировать объекты по убыванию $\hat{p}_i$
    \State Для каждого объекта в порядке убывания вероятности:
    \State \quad Рассмотреть его вероятность как порог $\tau$.
    \State \quad Вычислить Precision и Recall при этом пороге.
    \State Вычислить площадь под ломаной кривой (правило трапеций)
\end{algorithmic}

\underline{2. Формула через ранги}

Для бинарной классификации PR AUC можно вычислить как:
\[ \text{PR AUC} = \frac{\sum_{i=1}^{n_+} \frac{i}{r_i}}{n_+} \]
где:
\begin{itemize}
    \item $n_+$ --- количество положительных объектов
    \item $r_i$ --- ранг $i$-го положительного объекта после сортировки по убыванию вероятности
\end{itemize}

\underline{Интерпретация PR AUC}

\begin{itemize}
    \item \textbf{PR AUC = 1:} идеальный классификатор
    \item \textbf{PR AUC = положительная доля:} случайный классификатор
    \item \textbf{PR AUC близко к 0:} классификатор хуже случайного
\end{itemize}

\textbf{Важно:} Базовый уровень (случайный классификатор) для PR кривой --- горизонтальная линия на уровне положительной доли:
\[ \text{Precision}_{\text{random}} = \frac{n_+}{n} = P \]

\underline{Средние метрики для многоклассовой классификации}

\begin{enumerate}
    \item \textbf{Macro-average PR AUC:}
    \[ \text{Macro PR AUC} = \frac{1}{C} \sum_{c=1}^C \text{PR AUC}_c \]
    Здесь $\text{ PR AUC}_c$ вычисляется для класса $c$ против всех остальных
    
    \item \textbf{Micro-average PR AUC:} Объединение предсказаний для всех классов и вычисление общей PR кривой
    
    \item \textbf{Weighted-average PR AUC:}
    \[ \text{Weighted PR AUC} = \sum_{c=1}^C w_c \cdot \text{PR AUC}_c, \quad w_c = \frac{N_c}{N} \]
\end{enumerate}

\underline{Преимущества PR AUC для несбалансированных данных}

\begin{enumerate}
    \item \textbf{Чувствительность к изменениям в миноритарном классе:} небольшие изменения в предсказаниях
    миноритарного класса сильно влияют на PR кривую
    \item \textbf{Наглядность trade-off:} ясно показывает компромисс между точностью и полнотой
    \item \textbf{Информативность при high-recall:} важна в задачах, где нужно найти как можно больше положительных
    примеров (медицина, безопасность)
    \item \textbf{Устойчивость к изменению распределения:} PR кривая зависит только от ранжирования, а не от абсолютных
    значений вероятностей
\end{enumerate}

\subsection{Многоклассовые расширения}

Если классов становится больше двух, расчёт метрик усложняется.
Если задача классификации на $K$ классов ставится как $K$ задач об отделении класса $i$ от остальных
$(i=1, \dots, K)$, то для каждой из них можно посчитать свою матрицу ошибок.
Затем есть два варианта получения итогового значения метрики из $K$ матриц ошибок:

\subsubsection{Микро- и макро-усреднение}

\begin{enumerate}
    \item Усредняем элементы матрицы ошибок (TP, FP, TN, FN) между бинарными классификаторами, например
    $TP = \frac{1}{K} \sum_{i=1}^K TP_i$.
    Затем по одной усреднённой матрице ошибок считаем Precision, Recall, F-меру.
    Это называют микроусреднением.
    \item Считаем Precision, Recall для каждого классификатора отдельно, а потом усредняем.
    Это называют макроусреднением.
\end{enumerate}
Порядок усреднения влияет на результат в случае дисбаланса классов.
Показатели TP, FP, FN -- это счётчики объектов.
Пусть некоторый класс обладает маленькой мощностью (обозначим её MM).
Тогда значения TP и FN при классификации этого класса против остальных будут не больше MM, то есть тоже маленькие.
Про FP мы ничего уверенно сказать не можем, но скорее всего при дисбалансе классов классификатор не будет
предсказывать редкий класс слишком часто, потому что есть большая вероятность ошибиться.
Так что FP тоже мало.
Поэтому усреднение первым способом сделает вклад маленького класса в общую метрику незаметным.
А при усреднении вторым способом среднее считается уже для нормированных величин, так что вклад каждого класса будет
одинаковым.

Рассмотрим пример.
Пусть есть датасет из объектов трёх цветов: жёлтого, зелёного и синего.
Жёлтого и зелёного цветов поровну -- 21 и 20 объектов соответственно, а синих объектов всего 4.

Модель по очереди для каждого цвета пытается отделить объекты этого цвета от объектов оставшихся двух цветов.
Результаты классификации проиллюстрированы матрицей ошибок.
Модель «покрасила» в жёлтый 25 объектов, 20 из которых были действительно жёлтыми (левый столбец матрицы).
В синий был «покрашен» только один объект, который на самом деле жёлтый (средний столбец матрицы).
В зелёный -- 19 объектов, все на самом деле зелёные (правый столбец матрицы).
Макроусреднение лучше отражает тот факт, что синий цвет, которого в датасете было совсем мало, модель практически
игнорирует.